\documentclass[11pt]{article}
\usepackage[margin=1in]{geometry}
\usepackage{amsmath,amssymb}
\usepackage{graphicx}
\usepackage{booktabs}
\usepackage{array}
\usepackage[hidelinks]{hyperref}
\usepackage{caption}
\newcommand{\numu}{\nu_\mu}
\newcommand{\cmsq}{\mathrm{cm}^2}
\newcommand{\FIG}{/home/t2k/nowak/MicroBooNE/working_xsec_analyzer/xsec_analyzer/unfold_output}
\graphicspath{{\FIG/}}
\title{\bf CC1$\mu$1$\pi$Xp Selection and Extraction Procedure\\[2pt]
\large MicroBooNE NuMI Run 1 FHC}
\author{Internal technical note --- companion to \texttt{NuMIInternalNoteCC1pi}}
\date{Draft}

\begin{document}
\maketitle

\begin{abstract}
\noindent
This note documents, end to end, the selection and cross-section extraction for
the $\numu$ charged-current single-charged-pion measurement on argon
(\texttt{CC1mu1piXp}) with MicroBooNE Run~1 FHC NuMI data. It covers the signal
definition, the samples, the full cut sequence and particle identification, the
measured selection performance, and the extraction chain (background
subtraction, response matrix, Wiener--SVD unfolding with the additional smearing
matrix $A_C$, and the flux/target normalisation). All numbers are read from the
code, configs, or the universe files; where a quantity is currently configured
as a fake-data study rather than a data measurement this is stated explicitly.
\end{abstract}

\section{Signal definition and phase space}
An event is signal if, at truth level, the neutrino vertex is inside the
fiducial volume (FV), the interaction is charged current, the neutrino is
$\numu$ ($|\mathrm{mc\_nu\_pdg}|=14$, i.e.\ $\numu+\bar\numu$), and the final
state has exactly one muon, exactly one charged pion, no neutral pions, no
kaons, and no heavier mesons; any number of protons is allowed (the ``Xp'').
The measured phase space is
\begin{equation}
p_\mu > 0.15~\mathrm{GeV}/c,\qquad
p_\pi > 0.175~\mathrm{GeV}/c,\qquad
\theta_{\mu\pi} < 2.6~\mathrm{rad}.
\end{equation}
The same phase-space cuts are applied in \texttt{define\_signal()} and
\texttt{categorize\_event()}, so signal and category assignment cannot disagree.

\begin{table}[htbp]\centering
\caption{Fiducial volume and target count. The dead $z$ region is excluded.}
\begin{tabular}{lll}
\toprule
Quantity & Value & Notes\\
\midrule
$x$ range & $10 < x < 246$ cm & \\
$y$ range & $-101 < y < 101$ cm & \\
$z$ range & $10 < z < 986$ cm & dead region $675.1<z<775.1$ excluded\\
Net FV volume & $4.176\times10^{7}~\cmsq\cdot$cm & \\
LAr density & $1.3836$ g/cm$^3$ & \\
$N_{\mathrm{Ar}}$ (targets) & $8.710\times10^{29}$ nuclei & per-nucleus normalisation\\
\bottomrule
\end{tabular}
\end{table}

\section{Samples}
Run~1 FHC, from \texttt{configs/file\_properties\_numi.txt}.
\begin{table}[htbp]\centering
\caption{Analysis samples. The \texttt{onBNB} slot currently holds GENIE fake
data (the detVar central-value sample) for a closure study; the real beam-on
sample is commented out.}
\begin{tabular}{lll}
\toprule
Role & Sample & Purpose\\
\midrule
\texttt{onBNB} & GENIE detVar-CV (fake data) & pseudo-data (closure)\\
\texttt{extBNB} & beam-off (Run 1+3) & cosmic/EXT background, $3\,821\,593$ triggers\\
\texttt{numuMC} & NuMI overlay CV (GENIE) & response, efficiency, background model\\
\texttt{dirtMC} & dirt overlay & out-of-cryostat background\\
\texttt{detVarCV}$+$8 & detector variations & detector systematics\\
\bottomrule
\end{tabular}
\end{table}

\section{Selection}
\subsection{Cut sequence}
The cuts are applied in order; an event passes only if all are satisfied.
\begin{table}[htbp]\centering
\caption{The ten selection cuts (\texttt{CC1mu1piXp} selection class).}
\begin{tabular}{clp{8.2cm}}
\toprule
\# & Cut & Description\\
\midrule
1 & software trigger & beam/software trigger\\
2 & vertex in FV & reconstructed neutrino vertex inside the FV\\
3 & topological score & Pandora topological score (neutrino-like)\\
4 & muon track-like & muon candidate is track-like\\
5 & pion contained & pion candidate is contained\\
6 & muon gap & muon candidate not in a detector gap\\
7 & pion gap & pion candidate not in a detector gap\\
8 & shower cut & shower-multiplicity requirement\\
9 & opening angle & $\theta_{\mu\pi}<2.6$ rad\\
10 & final multiplicity & final track-multiplicity requirement\\
\bottomrule
\end{tabular}
\end{table}

\subsection{Particle identification}
Three boosted decision trees (BDTs) are applied via \texttt{TMVA::Reader}.
\begin{table}[htbp]\centering
\caption{PID BDTs. Inputs: three-plane Bragg likelihood ratios, LLR PID score,
track score, track length, and SCE-corrected track end position (the pion
reader omits track length).}
\begin{tabular}{lll}
\toprule
Reader & Weights & Role\\
\midrule
\texttt{tmvaReader} & \texttt{dataset\_MIP\_BDT\_no\_len} & MIP identification\\
\texttt{tmvaReader\_mu} & \texttt{dataset\_muon\_BDT} & muon --- \emph{booked but not evaluated}\\
\texttt{tmvaReader\_pi} & \texttt{dataset\_pion\_BDT\_no\_len} & pion\\
\bottomrule
\end{tabular}
\end{table}

\subsection{Selection performance}
Measured on the \texttt{numuMC} overlay (CV-weighted) for the cos$\theta_\mu$
binning; the integrated efficiency and purity are properties of the selection
and are essentially observable-independent.
\begin{table}[htbp]\centering
\caption{Integrated selection performance (CV-weighted \texttt{numuMC}).}
\begin{tabular}{ll}
\toprule
Quantity & Value\\
\midrule
Total selected (reco) & $\approx 3.71\times10^{3}$ events\\
Selected signal & $\approx 2.21\times10^{3}$ events\\
All true signal (denominator) & $\approx 1.42\times10^{4}$ events\\
\textbf{Purity} & \textbf{0.60}\\
\textbf{Efficiency} & \textbf{0.16}\\
\bottomrule
\end{tabular}
\end{table}

\noindent Figure~\ref{fig:reco} shows the selected reconstructed
cos$\theta_\mu$ spectrum split into MC signal and MC$+$EXT background, and the
selection efficiency versus true cos$\theta_\mu$.
\begin{figure}[htbp]\centering
\includegraphics[width=0.62\textwidth]{selection_costhmu_step1_reco_spectrum.pdf}\\[4pt]
\includegraphics[width=0.62\textwidth]{selection_costhmu_step4_efficiency.pdf}
\caption{Top: selected reco cos$\theta_\mu$ spectrum, data (fake data) vs MC
signal $+$ MC/EXT background. Bottom: selection efficiency vs true
cos$\theta_\mu$. The same two stages are produced for every observable.}
\label{fig:reco}
\end{figure}

\section{Binning}
\begin{table}[htbp]\centering
\caption{Differential observables and binning. Each true-bin count is one
greater than the reco count; the extra true bin is the background
(non-signal) bin.}
\begin{tabular}{lcl}
\toprule
Observable & Bins (true/reco) & Range\\
\midrule
$p_\mu$ & 23/22 & $0.15$--$3.00$ GeV/$c$\\
$p_\pi$ & 6/5 & $0.175$--$1.00$ GeV/$c$\\
$\cos\theta_\mu$ & 13/12 & $-1$ to $1$\\
$\cos\theta_\pi$ & 13/12 & $-1$ to $1$\\
$\theta_{\mu\pi}$ & 10/9 & $0$--$2.6$ rad\\
\bottomrule
\end{tabular}
\end{table}

\section{Cross-section extraction}
\subsection{Formula}
The flux-integrated differential cross section in true bin $\alpha$ is
\begin{equation}
\left.\frac{d\sigma}{dx}\right|_\alpha =
\frac{1}{N_{\mathrm{Ar}}\,\Phi\,\Delta x_\alpha}
\sum_a U_{\alpha a}\,(D_a - B_a),
\end{equation}
where $D_a$ is the measured reco rate in bin $a$, $B_a$ the estimated
background (EXT $+$ MC), $U_{\alpha a}$ the unfolding matrix (absorbing the
reco$\to$true mapping and the efficiency), $N_{\mathrm{Ar}}$ the FV target
count, $\Phi$ the integrated flux, and $\Delta x_\alpha$ the bin width. Results
are per argon nucleus, in $10^{-38}~\cmsq$ differential in the observable.

\subsection{Background subtraction and response}
The measured reco data ($D_a$, the \texttt{onBNB} histogram) has the EXT$+$MC
background subtracted, then is unfolded with the response built from the
\texttt{numuMC} overlay. Figure~\ref{fig:extract} shows the background-subtracted
spectrum (unfolding input) and the smearceptance (response) matrix.
\begin{figure}[htbp]\centering
\includegraphics[width=0.60\textwidth]{selection_costhmu_step2_bkgd_subtraction.pdf}\\[4pt]
\includegraphics[width=0.60\textwidth]{selection_costhmu_step3_smearing_matrix.pdf}
\caption{Top: background-subtracted reco cos$\theta_\mu$ with the MC signal
overlaid (the unfolding input). Bottom: the reco-vs-true smearceptance
(response) matrix inverted by the Wiener--SVD unfolding.}
\label{fig:extract}
\end{figure}

\subsection{Unfolding and the additional smearing matrix}
Unfolding uses Wiener--SVD with the filter enabled and second-derivative
regularisation (\texttt{Unfold WienerSVD 1 second-deriv}), chosen over
D'Agostini after a variant study because it suppresses the bin-to-bin
oscillation from the sparse Run~1 statistics. The Wiener--SVD result carries an
\emph{additional smearing matrix} $A_C$: the unfolded estimator is
$\hat{x}=A_C\,x_{\mathrm{true}}$, so any model or truth distribution must be
smeared by $A_C$ before it is compared to the unfolded data. $A_C$ does not in
general conserve the total, so the unfolded total cross section differs from the
raw truth total by the $A_C$ normalisation; the correct closure comparison is
$\hat{x}$ versus $A_C\,x_{\mathrm{true}}$.

\subsection{Normalisation and the conversion factor}
The conversion from unfolded true signal events to cross section is
$\mathrm{conv} = N_{\mathrm{Ar}}\,\Phi/10^{-38}$, with the integrated flux
$\Phi = \mathrm{POT}\times\phi$ and the authoritative NuMI FHC flux constant
$\phi = 6.81159\times10^{-10}~\numu/\mathrm{cm}^2/\mathrm{POT}$ for $E>60$ MeV
($= 4.43515\times10^{-10}$ [$\numu$] $+ 2.37644\times10^{-10}$ [$\bar\numu$]).
The flux is the active-volume-averaged flux used as an approximation for the FV
flux (a documented few-percent residual); no FV flux correction is applied.
\begin{table}[htbp]\centering
\caption{Conversion-factor inputs (fake-data POT).}
\begin{tabular}{ll}
\toprule
Quantity & Value\\
\midrule
$N_{\mathrm{Ar}}$ (FV) & $8.710\times10^{29}$\\
Flux constant $\phi$ ($E>60$ MeV) & $6.81159\times10^{-10}~\numu/\cmsq/$POT\\
POT & $7.630913\times10^{20}$\\
Integrated flux $\Phi$ & $5.198\times10^{11}~\numu/\cmsq$\\
$\mathrm{conv} = N_{\mathrm{Ar}}\Phi/10^{-38}$ & $4.527\times10^{3}$\\
\bottomrule
\end{tabular}
\end{table}

\section*{Status note}
The chain is currently configured as a fake-data study: the \texttt{onBNB} slot
holds GENIE fake data and the real beam-on sample is commented out. All
selection-performance and extraction numbers above are therefore validation
figures, not a data measurement. See \texttt{NuMIInternalNoteCC1pi} for the
fake-data closure results and the four-generator comparison.

\end{document}
